\section{Supuestos del modelo jerárquico y su verificación}\label{app:supuestos}

La tabla~\ref{tab:supuestos} confronta los 13 supuestos del modelo de la
ecuación~\eqref{eq:jerarquico} (M6) con la evidencia disponible: uno queda
verificado, siete tienen evidencia parcial,
cuatro carecen de contraste y uno no se cumple. \emph{Parcial} significa que la evidencia es
indirecta o cubre solo una parte del supuesto; \emph{sin contraste}, que no se
ejecutó ninguna prueba. Las cifras de M6 son salidas de
\texttt{docs/modelo\_jerarquico.pdf}, el render de
\texttt{notebooks/modelo\_jerarquico.qmd}; el modelo no se reajustó para este
apéndice. La fila de la pertenencia municipal es un cálculo directo sobre
\texttt{espacial.celdas()} y \texttt{atus\_zmm\_limpio\_geo.parquet}, ausente de
los \emph{notebooks}.

\begin{table*}
	\centering\scriptsize
	\caption{Supuestos del modelo jerárquico, verificación aplicada, resultado y
		estado. Las desviaciones posteriores (d.\,p.), $\hat{R}$ y el tamaño efectivo
		de muestra (ESS) provienen de ArviZ; \texttt{az.summary} redondea $\hat{R}$ a dos decimales.}
	\label{tab:supuestos}
	\begin{tabular}{@{}>{\raggedright\arraybackslash}p{0.19\textwidth}>{\raggedright\arraybackslash}p{0.33\textwidth}>{\raggedright\arraybackslash}p{0.33\textwidth}>{\raggedright\arraybackslash}p{0.07\textwidth}@{}}
		\toprule
		Supuesto & Verificación & Resultado & Estado \\
		\midrule
		Verosimilitud NB con $\theta$ común, sin inflación de ceros
		& Poisson, NB2 y ZINB se compararon solo en el GLM celda$\times$mes
		(§\ref{sec:res-conteo}). En M6 hay cobertura y CRPS en 2024, sin rootograma,
		residuales cuantil ni conteo de ceros esperados
		& $\theta = 12.2$ (d.\,p.\ 0.447); cobertura del intervalo del 90\,\%: 87.2\,\%;
		25.8\,\% de ceros en el panel de entrenamiento, sin esperado que comparar
		& Parcial \\
		\addlinespace
		Independencia condicional de $Y_{ia}$ dado $\mu_{ia}$
		& Sin contraste: el $I$ de Moran del residual se calculó para NB2, no para M6
		& NB2 en 2024: 0.303, y 0.056 sin la media municipal; M6: no calculado
		& Sin contraste \\
		\addlinespace
		Separabilidad espacio--tiempo dentro del año
		& Sin contraste: no hay prueba de interacción celda$\times$mes. Todos los GLM
		imponen un efecto de mes común, y la correlación entre años no mide la forma
		intraanual
		& ---
		& Sin contraste \\
		\addlinespace
		$u_i$ constante en 2019--2023
		& Spearman entre los conteos por celda de cada par de años
		(§\ref{sec:res-espacial}). M6 no tiene indicador de periodo
		& 0.833 entre 2019 y 2024; 0.863--0.896 dentro de 2019--2021, 0.884--0.906
		dentro de 2022--2024 y 0.800--0.839 entre bloques: el corte de 2022 queda sin
		modelar en $u_i$
		& Parcial \\
		\addlinespace
		Caminata aleatoria de $m_{j,a}$
		& Trayectorias a posteriori frente a las series observadas
		(figura~\ref{fig:res-trayectorias}): comparación visual, dentro de muestra y
		sin estadístico
		& Reproducen los saltos de los municipios con cambio de régimen, a los que se
		ajustaron
		& Parcial \\
		\addlinespace
		Innovaciones $t_3$
		& $\sigma_{rw}$ con $t_3$ frente a innovaciones normales. $\nu = 3$ se fijó y no
		se estimó; la variante normal no tiene código en el \texttt{.qmd}
		& 0.289 (d.\,p.\ 0.036) con $t_3$ y 0.61 con normales, esta sin intervalo
		& Parcial \\
		\addlinespace
		Persistencia $m_{j,2024} = m_{j,2023}$
		& Observado frente a predicho en 2024, en total y por municipio
		(§\ref{sec:res-validacion})
		& Total: 71,295 contra 70,810 (cociente predicho/observado 0.993). Por municipio,
		observado/predicho de 0.30 (Abasolo) a 2.56 (El Carmen): el total acierta porque
		errores de signo opuesto se compensan (§\ref{sec:disc-regimen})
		& No se cumple \\
		\addlinespace
		Grafo de vecindad completo
		& \texttt{espacial.celdas()} falla si las aristas de $A$ no coinciden con los
		pares de centroides a distancia $\sqrt{3}\,\ell$; el \emph{notebook} cuenta
		las componentes conexas
		& 5,325 aristas, 4.69 vecinas por celda. No conexo: 35 componentes, la mayor
		con 2,126 celdas (93.7\,\%). Ver el párrafo \emph{Grafo de vecindad}
		& Verificado \\
		\addlinespace
		Cada celda pertenece a un solo municipio (el modal)
		& Sin contraste en el \emph{notebook}. Cálculo directo: accidentes fuera del
		municipio modal de su celda
		& 10,559 de 379,160 accidentes (2.8\,\%) caen fuera; 430 de 2,270 celdas
		(18.9\,\%) tienen más de un municipio
		& Parcial \\
		\addlinespace
		Exogeneidad de $u_i$ respecto al municipio
		& Sin contraste: el \emph{notebook} cita la construcción (CAR centrado en cero),
		que no es una prueba, y $\beta_0$ compite con la media de $\mathbf{u}$
		& ---
		& Sin contraste \\
		\addlinespace
		$m_{j,a}$ mide reporte y no riesgo
		& No contrastable con ATUS, sin una fuente externa de exposición o de
		siniestralidad (§\ref{sec:disc-regimen})
		& ---
		& Sin contraste \\
		\addlinespace
		Convergencia del MCMC
		& $\hat{R} < 1.01$, ESS mínimo de 400 y cero divergencias en $\log\mu$
		(§\ref{sec:met-jerarquico}); en los hiperparámetros se reportan sin umbral
		& Cero divergencias. Cumplen $\log\mu$ ($\hat{R}$ máx.\ 1.002, ESS mín.\ 4,289),
		$\sigma_u$ y $\theta$ ($\hat{R}$ 1.00) y $\sigma_{rw}$ (1.01, ESS 259). Quedan en
		1.02 $\alpha$ (ESS 219), $\beta_0$ (ESS 149) y $\sigma_m$ (ESS 78)
		& Parcial \\
		\addlinespace
		Los priors no dictan el resultado
		& Reajuste con priors alternativos en $\alpha$, $\sigma_u$ y $\theta$
		(tabla~\ref{tab:priors}), 1,000 muestras por cadena. Sin alternativa para $\beta_0$,
		$\sigma_m$ ni $\sigma_{rw}$
		& $\alpha$ 0.919→0.921; $\theta$ 12.215→12.146; $\sigma_{rw}$ 0.289→0.287;
		$\sigma_u$ 2.606 sin cambio; CRPS 6.097→6.101; PAI 8.439 sin cambio. Dos
		divergencias; $\hat{R}$ máx.\ de $\log\mu$ 1.009
		& Parcial \\
		\bottomrule
	\end{tabular}
\end{table*}

\paragraph{Grafo de vecindad.} Una versión anterior de
\texttt{espacial.rejilla\_hexagonal()} calculaba los vértices de cada hexágono
por separado, con senos y cosenos, y los vértices compartidos diferían en unos
$4.7\times10^{-10}$\,m. \texttt{weights.Queen} solo une polígonos con
coordenadas idénticas, así que el grafo perdía 1,085 de 5,284 aristas
(20.5\,\%), quedaba en 120 componentes y contaba como islas a 29 celdas que sí
tenían una vecina con accidentes. Las aserciones de simetría y ausencia de
islas no lo detectaban. La rejilla actual calcula los vértices sobre una
retícula entera, de modo que dos hexágonos vecinos llegan a los mismos
flotantes, y \texttt{espacial.celdas()} contrasta el número de aristas con los
pares de centroides a distancia $\sqrt{3}\,\ell$, un criterio independiente de
la geometría de los polígonos. Con la corrección, la partición pasa de 2,241 a
2,270 celdas (74 islas en lugar de 103) y el $I$ de Moran global de 0.648 a
0.654. Con el mismo muestreo de 1,000 + 2,000 iteraciones, $\alpha$ pasa de
0.848 a 0.917, $\sigma_u$ de 2.38 a 2.61 y el CRPS de M6 de 6.166 a 6.093.

El grafo corregido sigue sin ser conexo: 35 componentes, y la mayor reúne el
93.7\,\% de las celdas. El modelo está bien definido, porque toda celda
conserva al menos una vecina y $D - \alpha A$ es definida positiva para
$\alpha < 1$ aunque el grafo esté desconexo. El encogimiento hacia las vecinas,
en cambio, solo opera dentro de cada componente.

\paragraph{Convergencia.} §\ref{sec:met-jerarquico} exige el umbral a
$\log\mu$ y lo reporta sin exigirlo en los hiperparámetros. $\alpha$,
$\beta_0$ y $\sigma_m$ quedan en $\hat{R} = 1.02$ con 3,000 iteraciones de
calentamiento y 4,000 muestras por cadena; con la mitad, $\alpha$ daba 1.04.
La mezcla lenta de $\alpha$ es estructural: cerca de $\alpha = 1$ el CAR propio
confunde la dependencia con $\sigma_u$. Su media es estable entre corridas y
priors (0.917--0.921). El \emph{notebook} no calcula $\hat{R}$ para $m_{j,a}$
ni para $u_i$ por separado: la convergencia de $\log\mu$ cubre su suma con
$\beta_0$, y la lectura individual de ambos efectos depende de que la mezcla
lenta de $\beta_0$ no se traslade a ellos.

\paragraph{Alcance de la evidencia.} La separabilidad y la exogeneidad figuran
como verificadas en la tabla de supuestos del \emph{notebook}, pero la evidencia
citada allí mide otra cosa: la correlación de Spearman mide estabilidad entre
años, no forma dentro del año, y un CAR centrado en cero es una elección del
modelo, no un resultado. La afirmación de que $u_i$ queda neto del municipio
(§\ref{sec:disc-regimen}) está condicionada a ambos supuestos y a que $m_{j,a}$
absorba todo el nivel. La validación en 2024 sostiene la distribución predictiva
en conjunto (§\ref{sec:res-validacion}), en una sola ventana
(§\ref{sec:disc-limitaciones}), y no permite atribuir el acierto a un supuesto
concreto.
